\documentclass[../../../include/open-logic-section]{subfiles}

\begin{document}

\olfileid{his}{set}{infinitesimals}
\olsection{અનંતસૂક્ષ્મો અને વિકલન}

ન્યૂટન અને લાઇબ્નિત્સે સત્તરમી સદીના અંતે સ્વતંત્ર રીતે
કલનશાસ્ત્ર શોધ્યું. કલનશાસ્ત્રનો એક ખાસ મહત્વનો ઉપયોગ
\emph{વિકલન} હતો. આશરે કહીએ તો, વિકલનનો હેતુ કોઈ
બિંદુએ વિધેયના “ફેરફારના દર” અથવા ઢાળની કલ્પના
આપવાનો છે.

આ વિચાર સમજાવવા એક ચિત્ર જુઓ. નીચે કાળા રંગે
દર્શાવેલું વિધેય $f(x) = \nicefrac{x^2}{4} + \nicefrac{1}{2}$
વિચારો:
\begin{center}
	\begin{tikzpicture}[scale=1]
	\draw[->, gray] (-1,0) -- (4.25,0) node[right] {$x$};
	\draw[->, gray] (0,-0.25) -- (0,5.25) node[above] {$f(x)$};
	
	\foreach \x/\xtext in {1/1, 2/2, 3/3, 4/4}
	\draw[shift={(\x,0)}, gray] (0pt,2pt) -- (0pt,-2pt) node[below] {$\xtext$};
	
	\foreach \y/\ytext in {1/1, 2/2, 3/3, 4/4, 5/5}
	\draw[shift={(0,\y)}, gray] (2pt,0pt) -- (-2pt,0pt) node[left] {$\ytext$};
	
	\draw[oldiagcolorC] (0.5, 9/16) -- (3.5, 57/16) -- (3.5, 9/16)--cycle;
	\draw[oldiagcolorD] (.5, 9/16) -- (2.5, 33/16) -- (2.5, 9/16) -- cycle;
	\draw[oldiagcolorE] (.5, 9/16) -- (1.5, 17/16) -- (1.5, 9/16) -- cycle;
	\draw[thick] (-1,.75) parabola bend (0,0.5) (4,4.5);
	\end{tikzpicture}
\end{center}
ધારો કે આપણે $c =
\nicefrac{1}{2}$ પર વિધેયનો ઢાળ શોધવો છે.
શરૂઆતમાં એવું ત્રિકોણ દોરીએ કે જેની કર્ણરેખાનો ઢાળ
આ બિંદુનો ઢાળ અંદાજે આપે; ઉપરનો લાલ ત્રિકોણ
એવું ઉદાહરણ છે. ત્રિકોણના પાયાની લંબાઈ
$\beta$ હોય, તો તેની ઊંચાઈ
$f(\nicefrac{1}{2}+\beta) - f(\nicefrac{1}{2})$
છે. તેથી કર્ણરેખાનો ઢાળ છે:
\[
\frac{f(\nicefrac{1}{2}+\beta) - f(\nicefrac{1}{2})}{\beta}.
\]
આથી પાયાની લંબાઈ~$3$ ધરાવતા આપણા
!!{colorC} ત્રિકોણનો ઢાળ બરાબર~$1$ છે.
તેના કરતાં નાના, પાયાની લંબાઈ~$2$
ધરાવતા !!{colorD} ત્રિકોણની કર્ણરેખા વધુ
સારો અંદાજ આપે છે; તેનો ઢાળ
$\nicefrac{3}{4}$ છે. પાયાની લંબાઈ~$1$
ધરાવતો હજી નાનો !!{colorE} ત્રિકોણ
વધુ સારો અંદાજ આપે છે; તેનો ઢાળ
$\nicefrac{1}{2}$ છે.

વધુ ને વધુ નાના ત્રિકોણોથી વધુ ને વધુ સારા અંદાજ
મળે છે. તેથી આપણે કદાચ કહી શકીએ કે
\emph{અનંતસૂક્ષ્મ} લંબાઈના પાયાવાળા ત્રિકોણની
કર્ણરેખા આપણને $c =
\nicefrac{1}{2}$ પરનો ઢાળ પોતે આપે છે.
આ રીતે બિંદુ $c$ પર વિધેય $f$ના
(પ્રથમ) વિકલનફળનું સૂત્ર મળે:
\[
{f'}(c) = \frac{f(c+\beta) - f(c)}{\beta} \text{ જ્યાં $\beta$ અનંતસૂક્ષ્મ છે.}
\]
આશરે ન્યૂટન અને લાઇબ્નિત્સે આવું જ કહ્યું હતું.

પરંતુ એમ કહેતાં તેમને પૂછવું પડે: \emph{અનંતસૂક્ષ્મ}
એટલે શું? અહીં ગંભીર મૂંઝવણ ઊભી થાય છે.
જો $\beta = 0$ હોય, તો $0$થી ભાગાકાર થતો
હોવાથી $f'$ સુવ્યાખ્યાયિત નથી.
પરંતુ જો $\beta >
0$ હોય, તો આપણને ઢાળ પોતે નહીં,
ફક્ત તેનો \emph{અંદાજ} મળે છે.

આ ચિંતા પાછળથી લાદેલી નથી.
ન્યૂટનના અનુયાયીઓની ટીકા કરતાં બર્કલીને સાંભળો:
\begin{quote}
ચિહ્નો વડે કશુંક અથવા કશું જ નહીં દર્શાવી શકાય,
એ હું સ્વીકારું છું. તેથી મૂળ સંકેતન
$c + \beta$માં $\beta$ વધારો અથવા કશું જ નહીં
દર્શાવતું હોઈ શકે. પરંતુ તેનો જે અર્થ લો,
દલીલમાં એ જ અર્થ સાથે સુસંગત રહેવું જોઈએ;
બેવડા અર્થ પર આધાર રાખીને આગળ ન વધવું જોઈએ.
એમ કરવું સ્પષ્ટ કુતર્ક છે.
(\citealt[\S{}XIII]{Berkeley1734};
અગાઉના લખાણ સાથે મેળ કરવા ચલો બદલ્યાં છે)
\end{quote}
બર્કલી સામે અનંતસૂક્ષ્મ કલનનો બચાવ કરતાં
કોઈ કહી શકે કે “અનંતસૂક્ષ્મો”ની વાત માત્ર
રૂપક છે: \emph{ખરેખર નાનું} ત્રિકોણ લઈએ
તો સ્પર્શરેખાનો \emph{પૂરતો સારો} અંદાજ
મળી જાય. બર્કલી પાસે તેનો પણ જવાબ હતો:
ઇજનેરી માટે આ પૂરતું હોય, પરંતુ ગણિતના
\emph{દરજાને} તે હાનિ પહોંચાડે છે, કારણ કે
\begin{quote}
આપણને કહેવામાં આવે છે કે
\emph{in rebus mathematicis errores qu\`{a}m minimi
non sunt contemnendi}. [ગણિતમાં સૌથી નાની
ભૂલોની પણ અવગણના કરવી નહીં.]
\citep[\S{}IX]{Berkeley1734}
\end{quote}
ત્રાંસા અક્ષરોમાંનો ભાગ ન્યૂટનની પોતાની
\emph{Quadrature of Curves} (1704)
માંથી લગભગ શબ્દશઃ લીધેલો છે.

બર્કલીના તત્વજ્ઞાનાત્મક વાંધા ખૂબ
તીક્ષ્ણ છે. તેમ છતાં કલનશાસ્ત્ર
અસાધારણ સફળ રહ્યું, અને ગણિતશાસ્ત્રીઓ
ભૂલમાં પડ્યા વિના તેનો ઉપયોગ કરતા રહ્યા.

\end{document}
